% block: condition-number-heading page:17 type:section
\subsection{\BT{Condition Number}{条件数}}\label{sec:condition-number}
% block: perturbation-equation page:17 type:ordinary-paragraph
\BT{Consider $A\bm x=\bm b$; if the matrix $A$ and right-hand side vector $\bm b$ are slightly changed, then by how much will the solution be changed? To answer this question, let $\Delta A$ and $\Delta\bm b$ be perturbations of $A$ and $\bm b$. Then the change of the solution denoted by $\Delta\bm x$ satisfies the following relation:}{考虑 $A\bm x=\bm b$; 如果矩阵 $A$ 和右端向量 $\bm b$ 发生微小变化, 解会改变多少? 为回答这一问题, 设 $\Delta A$ 和 $\Delta\bm b$ 分别为 $A$ 和 $\bm b$ 的扰动. 以 $\Delta\bm x$ 表示的解的变化满足下列关系:}
\[
(A+\Delta A)(\bm x+\Delta\bm x)=\bm b+\Delta\bm b.
\]
\BT{Using $A\bm x=\bm b$ gives}{利用 $A\bm x=\bm b$ 得到}
\[
\Delta\bm x=A^{-1}[-\Delta A(\bm x+\Delta\bm x)+\Delta\bm b].
\]
\BT{Using Theorem~\ref{thm:matrix-norm-properties}-(\ref{part:matrix-norm-action}) and the triangle inequality (\ref{cond:vector-norm-triangle}), the $p$-norm of $\bm x$ can be evaluated by}{利用 Theorem~\ref{thm:matrix-norm-properties}-(\ref{part:matrix-norm-action}) 和三角不等式 (\ref{cond:vector-norm-triangle}), 可由下式估计 $\bm x$ 的 $p$-范数:}
\[
\|\Delta\bm x\|_p\le\|A^{-1}\|_p(\|\Delta A\|_p\|\bm x+\Delta\bm x\|_p+\|\Delta\bm b\|_p).
\]
\BT{Thus, we have}{因此, 有}
\begin{equation}\label{eq:relative-perturbation-bound}
\frac{\|\Delta\bm x\|_p}{\|\bm x+\Delta\bm x\|_p}
\le\|A\|_p\|A^{-1}\|_p
\left(\frac{\|\Delta A\|_p}{\|A\|_p}+
\frac{\|\Delta\bm b\|_p}{\|A\|_p\|\bm x+\Delta\bm x\|_p}\right).
\end{equation}
% block: condition-number-definition page:17 type:ordinary-paragraph
\BT{Here,}{这里,}
\begin{equation}\label{eq:p-norm-condition-number}
\kappa_p(A):=\|A\|_p\|A^{-1}\|_p
\end{equation}
\BT{is called the $p$-norm condition number. Inequality \eqref{eq:relative-perturbation-bound} implies that if the condition number is large, the computed approximate solution may be far from the exact solution. Note that this fact does not depend on what algorithms are used for solving linear systems.}{称为 $p$-范数条件数. 不等式 \eqref{eq:relative-perturbation-bound} 表明, 如果条件数很大, 计算得到的近似解可能远离精确解. 注意, 这一事实与求解线性方程组所使用的算法无关.}
% block: ill-conditioning page:17 type:ordinary-paragraph
\BT{If the condition number is large, then the linear systems are called ill-conditioned. We will see in Chapter~\ref{chap:krylov-classification} that the condition number also affects the speed of convergence of Krylov subspace methods.}{如果条件数很大, 则称该线性方程组是病态的. 我们将在 Chapter~\ref{chap:krylov-classification} 看到, 条件数也会影响 Krylov 子空间方法的收敛速度.}
\BT{In this book, we often use a 2-norm condition number $\kappa_2(A)$ that will be simply denoted by $\kappa$ or $\operatorname{cond}(A)$.}{本书经常使用 2-范数条件数 $\kappa_2(A)$, 并将其简记为 $\kappa$ 或 $\operatorname{cond}(A)$.}
